% TOP_PROBLEM: {"id": 2867, "title": "Kirby Problem 3.69", "category_id": 7, "category": {"id": 7, "name": "topology", "display_name": "Topology", "description": "Properties preserved under continuous deformations.", "slug": "topology", "order_index": 7, "created_at": "2026-07-31T15:26:25.671Z"}, "status": "open"}
% FIRST_POSTED: null
% ENTRY_KIND: statement_only
% Imported from: batch12.tex; UnsolvedMath ID 2867
\documentclass[11pt]{article}
\usepackage[margin=25mm]{geometry}
\usepackage{fontspec}
\setmainfont{Latin Modern Roman}
\usepackage{amsmath,amssymb,mathtools}
\usepackage{unicode-math}
\setmathfont{Latin Modern Math}
\usepackage{xurl,hyperref}
\hypersetup{colorlinks=true,urlcolor=blue}
\setcounter{secnumdepth}{0}
\setlength{\parindent}{0pt}
\setlength{\parskip}{5pt}
\setlength{\emergencystretch}{3em}
\newcommand{\sref}[2]{\hyperlink{source:#1:#2}{[S#2]}}
\newcommand{\cataloguescope}[1]{\paragraph{Recorded target.}#1}
\begin{document}
% UnsolvedMath ID 2867; source catalogue code KP-3.69
\setcounter{section}{2866}
\section{Kirby Problem 3.69}\label{2867}
{\small\sffamily UnsolvedMath ID 2867 · Topology}
\subsection{Definitions and mathematical statement}

\paragraph{Shared notation.}$\mathbb N=\{1,2,\ldots\}$, and $\mathbb Z,\mathbb Q,\mathbb R,\mathbb C$ denote the integers, rationals, reals and complex numbers. Any other conventions are specified locally.


An integral homology $3$-sphere is a closed oriented smooth $3$-manifold $Y$ with $H_*(Y;\mathbb Z)\cong H_*(S^3;\mathbb Z)$. Two such spheres $Y_0,Y_1$ are integrally homology cobordant if there is a compact connected oriented smooth $4$-manifold $W$ with $\partial W=-Y_0\sqcup Y_1$ for which both inclusions induce isomorphisms on integral homology. Their classes form the abelian group $\Theta^3_{\mathbb Z}$ under connected sum, with zero $[S^3]$ and inverse $[-Y]$.
The Rokhlin homomorphism $\mu:\Theta^3_{\mathbb Z}\to\mathbb Z/2\mathbb Z$ is $\mu([Y])=\sigma(X)/8\pmod2$, where $X$ is any smooth compact spin $4$-manifold bounded by $Y$ and $\sigma$ is its intersection-form signature. A nontrivial torsion element means $x\ne0$ with $mx=0$ for some $m\ge2$. For $n\ge1$, $nz$ means the sum of $n$ copies of $z$.
\paragraph{Statement recorded in the supplied source.}
Resolve all four questions:
\begin{enumerate}
\item Determine the isomorphism type of $\Theta^3_{\mathbb Z}$.
\item Does $\Theta^3_{\mathbb Z}$ contain a nontrivial torsion element?
\item Does it contain a torsion element $x$ with $\mu(x)=1$?
\item Does there exist $x=[Y]$ such that the set
\[
\{n\in\mathbb N:\ \exists z_n\in\Theta^3_{\mathbb Z},\ x=nz_n\}
\]
is infinite?
\end{enumerate}
Part (d) is the literal infinitely-many-$n$ condition in the source. It permits $x=0$, which satisfies the condition by taking $z_n=0$ for every $n$. Requiring a nonzero class, requiring divisibility by every $n$, or passing to the quotient by torsion are separate stronger formulations. \sref{2867}{2}
\subsection{Short English statement}
Determine the integral homology-cobordism group, whether it has nontrivial torsion or torsion with Rokhlin invariant one, and its stated infinitely-many-divisors condition. The literal last condition is already satisfied by the zero class.
\subsection{Sources}
\begin{description}
\item[\hypertarget{source:2867:1}{S1}] ULAM.AI and the UnsolvedMath Contributors, \emph{UnsolvedMath: A Curated Collection of Open Mathematics Problems}, record 2867 (KP-3.69). CC BY 4.0. \url{https://huggingface.co/datasets/ulamai/UnsolvedMath}.
\item[\hypertarget{source:2867:2}{S2}] R. İnanç Baykur, Robion C. Kirby and Daniel Ruberman (editors), \emph{K3: A New Problem List in Low-Dimensional Topology}, Mathematical Surveys and Monographs 295 (2026), Problem 3.69, p. 181 and following remarks; original author version. \url{https://math.berkeley.edu/sites/default/files/surv-295-ruberman-watermarked-author-pdf.pdf}.
\end{description}
\label{end:2867}
\end{document}
